\documentclass[a4,11pt]{wmecai}

\usepackage[brazil]{babel}
\usepackage[utf8]{inputenc}

%%%%%%%%%%%%%%%%%%%%%%%%%%%%%%%%%%%%%%%%%%%%%%%%%%%%%%%%%%%%%%
%% POR FAVOR, NÃO FAÇA MUDANÇAS NESSE PADRÃO QUE ACARRETEM  EM
%% ALTERAÇÃO NA FORMATAÇÃO FINAL DO TEXTO
%%%%%%%%%%%%%%%%%%%%%%%%%%%%%%%%%%%%%%%%%%%%%%%%%%%%%%%%%%%%%%

\usepackage{graphicx}
\usepackage{subfigure}
\usepackage[centertags]{amsmath}
\usepackage{indentfirst,amsmath,amsfonts,amssymb,amsthm,newlfont}
\usepackage{longtable}
\usepackage{cite}
\usepackage[usenames]{color}
\usepackage{booktabs}

\begin{document}

\title{Bankai Boruta: escalando seleção de features para modelos de classificação com robustez estatística}

\author{
{\large Nome do(a) Autor(a)}\thanks{email@exemplo.com}\\
{\small Instituição do(a) autor(a)}\\
}

\criartitulo

\section{Introdução}
A seleção de atributos é uma etapa importante em problemas de classificação: pode reduzir custo computacional, facilitar a interpretação e ajudar a identificar variáveis relevantes. Diferentemente de métodos que procuram apenas um subconjunto mínimo, o Boruta busca encontrar todos os atributos associados à resposta, comparando iterativamente a importância dos atributos observados com a de cópias aleatorizadas, chamadas atributos-sombra \cite{Kursa}. A repetição desse procedimento com florestas aleatórias pode, contudo, tornar-se onerosa em conjuntos de dados grandes, sobretudo quando há muitas variáveis contínuas.

Este trabalho apresenta o \textit{Bankai Boruta}, uma implementação em Rust do algoritmo Boruta voltada a classificação e integrada ao ecossistema Python e scikit-learn. A proposta explora duas frentes: (i) um caminho nativo, apoiado por uma floresta aleatória em Rust com cálculo nativo de importância; e (ii) a possibilidade de fornecer outros estimadores compatíveis, preservando o fluxo de uso familiar a profissionais de aprendizado de máquina. Para o caminho com o Bankai Random Forest, a discretização de variáveis contínuas em bins limita os pontos candidatos a divisão, estratégia relacionada às otimizações por histogramas empregadas em árvores de decisão eficientes, como no LightGBM \cite{Ke}.

A hipótese de trabalho é que a combinação da implementação nativa e da discretização reduza o tempo de seleção em bases maiores, mantendo decisões de seleção comparáveis às de uma implementação de referência. Essa hipótese deve ser avaliada empiricamente; este rascunho descreve o método e o protocolo, sem atribuir ganhos não medidos. As contribuições pretendidas são: uma implementação de alto desempenho do fluxo Boruta; suporte a diferentes provedores de importância, incluindo opções do Bankai Random Forest como ganho, contagem de divisões e SHAP; e uma avaliação controlada de custo e qualidade.

\section{Método}
\subsection{Seleção por atributos-sombra}
Seja $X\in\mathbb{R}^{n\times p}$ uma matriz de $n$ observações e $p$ atributos, e seja $y$ o vetor de classes. Em cada iteração $t$, o algoritmo permuta independentemente cada coluna ativa para criar atributos-sombra e concatena-os aos atributos reais. Um estimador ajustado a essa matriz fornece escores de importância. O maior escore entre os atributos-sombra constitui o limiar de comparação da iteração. Para cada atributo real, registra-se se sua importância excede esse limiar; decisões de aceitação ou rejeição são acumuladas ao longo das iterações e avaliadas estatisticamente. A repetição do processo permite distinguir atributos confirmados, rejeitados e ainda tentativos quando o limite de iterações é atingido \cite{Kursa}.

\subsection{Implementação e aceleração}
O controlador do Boruta, a geração de sombras, a poda de atributos rejeitados e as decisões estatísticas são implementados no núcleo Rust. No caminho nativo, o estimador de floresta e suas importâncias também são calculados em Rust, reduzindo a sobrecarga de chamadas entre linguagens. No caminho integrado, o Bankai Random Forest atua como provedor do estimador. A discretização de valores contínuos em um número configurável de bins permite reutilizar estatísticas agregadas por bin na busca de divisões, em vez de considerar cada valor distinto como candidato. Essa escolha oferece um compromisso entre custo e resolução: bins mais numerosos podem representar mais pontos de corte, enquanto menos bins tendem a reduzir o trabalho de treinamento.

A arquitetura também aceita estimadores classificadores compatíveis com scikit-learn que exponham um vetor de importância por atributo. Assim, o usuário pode trocar o modelo ou a medida de importância sem alterar o procedimento Boruta. Entre as opções do Bankai Random Forest estão importâncias por ganho, por divisão e SHAP; SHAP oferece uma abordagem unificada para atribuir importância às previsões de modelos \cite{Lundberg}. Essa flexibilidade permite que o caminho mais rápido e especializado conviva com fluxos baseados em outros modelos, com o custo adicional de materializar a matriz aumentada e ajustar o estimador em cada iteração.

\section{Protocolo de avaliação}
A avaliação deve comparar o Bankai Boruta com BorutaPy, mantendo iguais os dados, as sementes, os critérios estatísticos e, quando possível, o estimador e o número de árvores. Recomenda-se usar conjuntos sintéticos com atributos informativos, redundantes e irrelevantes, variando número de observações, número de atributos e proporção de valores contínuos, além de bases públicas de classificação. A implementação e as versões das dependências, as configurações de hardware e os parâmetros de cada execução devem ser registrados para permitir reprodução.

As métricas de desempenho incluem tempo total de seleção, pico de memória, número de iterações e, quando disponível, tempo gasto no ajuste do estimador. A qualidade da seleção em dados sintéticos pode ser quantificada por precisão, revocação e F1 em relação ao conjunto informativo conhecido. Em bases reais, nas quais a verdade de referência não é conhecida, devem ser reportadas a estabilidade da seleção entre sementes e o desempenho de um classificador treinado somente com os atributos selecionados, medido em partições de teste não usadas na seleção. A seleção deve ocorrer dentro de cada partição de treinamento para evitar vazamento de informação.

Um estudo de ablação deve comparar o caminho nativo, o caminho Bankai Random Forest com discretização e o mesmo estimador com diferentes resoluções de bins. Também deve comparar provedores de importância quando aplicável. Os resultados precisam apresentar medidas agregadas em execuções repetidas, por exemplo mediana e intervalo interquartil. Até a realização desses experimentos, não se afirma superioridade em velocidade ou qualidade: a discretização pode alterar as divisões da árvore e, consequentemente, as importâncias e decisões estatísticas.

\section{Reprodutibilidade e disponibilização}
O projeto é disponibilizado como software livre no GitHub e integra testes automatizados ao desenvolvimento, com esteira de integração contínua em múltiplas plataformas e verificações de dependências para reduzir riscos de dependências vulneráveis ou comprometidas. Esses mecanismos apoiam a manutenção e a reprodutibilidade da implementação, mas não substituem a descrição das versões, configurações e dados usados em cada experimento. O endereço do repositório, a licença e os identificadores das versões utilizadas devem ser incluídos na versão final da submissão.

\section{Conclusão}
O Bankai Boruta combina um núcleo de seleção em Rust, um caminho nativo com floresta aleatória e integração opcional com outros estimadores e medidas de importância. A discretização de variáveis contínuas no Bankai Random Forest constitui a principal hipótese de aceleração para bases de maior porte, enquanto a API compatível com scikit-learn preserva um fluxo de trabalho conhecido. O protocolo proposto permitirá medir separadamente custo computacional, qualidade e estabilidade da seleção. Como próximo passo, serão executados os experimentos descritos e incorporados seus resultados, incluindo limitações observadas e condições em que cada provedor é mais adequado.

\begin{thebibliography}{00}

\bibitem{Ke}
G. Ke, Q. Meng, T. Finley, T. Wang, W. Chen, W. Ma, Q. Ye, and T.-Y. Liu. LightGBM: A highly efficient gradient boosting decision tree. In \textit{Advances in Neural Information Processing Systems}, 30, 2017.

\bibitem{Kursa}
M. B. Kursa, A. Jankowski, and W. R. Rudnicki. Boruta---a system for feature selection. \textit{Fundamenta Informaticae}, 101(4):271--285, 2010. DOI: 10.3233/FI-2010-288.

\bibitem{Lundberg}
S. M. Lundberg and S.-I. Lee. A unified approach to interpreting model predictions. In \textit{Advances in Neural Information Processing Systems}, 30, 2017.

\end{thebibliography}
\end{document}
